\documentclass[tikz,border=6pt]{standalone}
\usepackage{mathdiagrams}
% scale_drawing_20m_10cm — drawing 10 cm ↔ actual 20 m
% Contract: no white boxes; no prose qualifiers ("on diagram", "in real life");
% measurements only — the formula carries the relationship.
% Check: 10 cm : 20 m = 10 : 2000 = 1 : 200.
\begin{document}
\begin{tikzpicture}[md]
  % diagram length
  \draw[md thick] (0,0) -- (5,0);
  \draw[md thin] (0,-0.15) -- (0,0.15);
  \draw[md thin] (5,-0.15) -- (5,0.15);
  \node[font=\sffamily\small, text=mdInk] at (2.5,0.32) {$10$ cm};

  % actual length
  \draw[md accent, thick] (0,-1.1) -- (5,-1.1);
  \draw[md thin] (0,-1.25) -- (0,-0.95);
  \draw[md thin] (5,-1.25) -- (5,-0.95);
  \node[font=\sffamily\small, text=mdInk] at (2.5,-1.48) {$20$ m};

  % scale formula (data)
  \node[font=\sffamily\small, text=mdInk] at (2.5,-2.25)
    {scale $= 10 \text{ cm} : 20 \text{ m} = 1 : 200$};
\end{tikzpicture}
\end{document}
